\documentclass[UTF8,a4paper]{ctexart}
\usepackage{amsmath,amssymb}
\usepackage{graphicx}
\usepackage{geometry}
\usepackage{enumitem}

\geometry{margin=2.0cm}
\graphicspath{{D:/Git/AmazingEducation/examples/case/bank/}}
\pagestyle{plain}
\begin{document}



\begin{center}
{\LARGE 高中数学综合练习（合成样卷）\par}
\medskip
{\large 案例题库}
\end{center}
\vspace{1em}

\section*{一、选择题}


\begin{enumerate}[leftmargin=2em,itemsep=0.8em]


\item 已知复数 $z=1+i$，则 $|z|$ 等于

\begin{enumerate}[label=\Alph*.,itemsep=0.2em,parsep=0pt]

\item $1$

\item $\sqrt{2}$

\item $2$

\item $\sqrt{3}$

\end{enumerate}


\par\noindent\textbf{【答案】}B
\par\noindent\textbf{【解析】}$|1+i|=\sqrt{1^2+1^2}=\sqrt{2}$。



\item 函数 $f(x)=x^2-4x+3$ 的最小值为

\begin{enumerate}[label=\Alph*.,itemsep=0.2em,parsep=0pt]

\item $-1$

\item $0$

\item $1$

\item $-2$

\end{enumerate}


\par\noindent\textbf{【答案】}A
\par\noindent\textbf{【解析】}$f(x)=(x-2)^2-1$，最小值为 $-1$。



\item 等差数列 $a_n$ 满足 $a_1=2, d=3$，则 $a_5=$

\begin{enumerate}[label=\Alph*.,itemsep=0.2em,parsep=0pt]

\item $11$

\item $14$

\item $12$

\item $10$

\end{enumerate}


\par\noindent\textbf{【答案】}B
\par\noindent\textbf{【解析】}$a_5=a_1+4d=2+12=14$。



\item 若 $\sin x=\frac{1}{2}$，且 $x\in(0,\pi)$，则 $x=$

\begin{enumerate}[label=\Alph*.,itemsep=0.2em,parsep=0pt]

\item $\frac{\pi}{6}$

\item $\frac{\pi}{3}$

\item $\frac{\pi}{2}$

\item $\frac{2\pi}{3}$

\end{enumerate}


\par\noindent\textbf{【答案】}A
\par\noindent\textbf{【解析】}在 $(0,\pi)$ 内满足 $\sin x=1/2$ 的解为 $x=\pi/6$ 或 $5\pi/6$，选项中只有 $\pi/6$。



\item 函数 $f(x)=e^x$ 的导数为

\begin{enumerate}[label=\Alph*.,itemsep=0.2em,parsep=0pt]

\item $xe^x$

\item $e^x$

\item $\ln x$

\item $x^{e}$

\end{enumerate}


\par\noindent\textbf{【答案】}B
\par\noindent\textbf{【解析】}指数函数 $e^x$ 的导数仍为 $e^x$。



\item 若向量 $\vec a=(1,2)$，$\vec b=(2,1)$，则 $\vec a \cdot \vec b=$

\begin{enumerate}[label=\Alph*.,itemsep=0.2em,parsep=0pt]

\item $3$

\item $4$

\item $5$

\item $2$

\end{enumerate}


\par\noindent\textbf{【答案】}B
\par\noindent\textbf{【解析】}$1\times2+2\times1=4$。



\item 集合 $A=\{1,2,3\}$，$B=\{2,3,4\}$，则 $A\cap B=$

\begin{enumerate}[label=\Alph*.,itemsep=0.2em,parsep=0pt]

\item $\{2,3\}$

\item $\{1,2\}$

\item $\{3,4\}$

\item $\{1,4\}$

\end{enumerate}


\par\noindent\textbf{【答案】}A
\par\noindent\textbf{【解析】}交集为两个集合共同元素。



\item 若随机变量 $X$ 服从二项分布 $B(2,\frac12)$，则 $P(X=0)=$

\begin{enumerate}[label=\Alph*.,itemsep=0.2em,parsep=0pt]

\item $\frac14$

\item $\frac12$

\item $\frac34$

\item $\frac18$

\end{enumerate}


\par\noindent\textbf{【答案】}A
\par\noindent\textbf{【解析】}$P(X=0)=(1/2)^2=1/4$。



\item 函数 $f(x)=\ln x$ 的定义域为

\begin{enumerate}[label=\Alph*.,itemsep=0.2em,parsep=0pt]

\item $(-\infty,\infty)$

\item $[0,\infty)$

\item $(0,\infty)$

\item $(-\infty,0)$

\end{enumerate}


\par\noindent\textbf{【答案】}C
\par\noindent\textbf{【解析】}对数函数要求 $x>0$。



\item 若 $a>0$，则 $\log_a 1=$

\begin{enumerate}[label=\Alph*.,itemsep=0.2em,parsep=0pt]

\item $0$

\item $1$

\item $a$

\item $-1$

\end{enumerate}


\par\noindent\textbf{【答案】}A
\par\noindent\textbf{【解析】}$a^0=1$。



\item 函数 $f(x)=x^3$ 在 $x=1$ 处的导数为

\begin{enumerate}[label=\Alph*.,itemsep=0.2em,parsep=0pt]

\item $1$

\item $2$

\item $3$

\item $6$

\end{enumerate}


\par\noindent\textbf{【答案】}C
\par\noindent\textbf{【解析】}$f'(x)=3x^2$，代入得 $3$。



\item 若等比数列 $a_n$ 满足 $a_1=2,q=2$，则 $a_4=$

\begin{enumerate}[label=\Alph*.,itemsep=0.2em,parsep=0pt]

\item $8$

\item $12$

\item $16$

\item $32$

\end{enumerate}


\par\noindent\textbf{【答案】}C
\par\noindent\textbf{【解析】}$a_4=a_1q^3=2\times8=16$。


\end{enumerate}
\vspace{1em}


\section*{二、填空题}


\begin{enumerate}[leftmargin=2em,itemsep=0.8em]


\item 计算：$\sum_{k=1}^{5} k = \underline{\qquad}$


\par\noindent\textbf{【答案】}$15$
\par\noindent\textbf{【解析】}$1+2+3+4+5=15$。



\item 函数 $f(x)=x^2$ 在 $x=3$ 处的导数为 $\underline{\qquad}$


\par\noindent\textbf{【答案】}$6$
\par\noindent\textbf{【解析】}$f'(x)=2x$，代入得 $6$。



\item 若 $\sin^2 x+\cos^2 x=$ $\underline{\qquad}$


\par\noindent\textbf{【答案】}$1$
\par\noindent\textbf{【解析】}三角恒等式。



\item 若 $a,b$ 为实数，$(a+b)^2-a^2-b^2=$ $\underline{\qquad}$


\par\noindent\textbf{【答案】}$2ab$
\par\noindent\textbf{【解析】}展开即可。


\end{enumerate}
\vspace{1em}


\section*{三、计算与证明题}


\begin{enumerate}[leftmargin=2em,itemsep=0.8em]


\item 已知函数 $f(x)=x^2-3x+2$。
(1) 求函数的零点；
(2) 求其在区间 $[0,3]$ 上的最小值。


\par\noindent\textbf{【答案】}(1) 解方程 $x^2-3x+2=0$，即 $(x-1)(x-2)=0$，故 $x=1$ 或 $x=2$。
(2) $f(x)=(x-\\frac{3}{2})^2-\\frac{1}{4}$，在 $[0,3]$ 上最小值为 $-\\frac{1}{4}$。

\par\noindent\textbf{【解析】}通过因式分解和配方法求解。



\item 已知等差数列 $a_n$ 满足 $a_1=3,d=2$。
(1) 求通项公式；
(2) 求前10项和。


\par\noindent\textbf{【答案】}(1) $a_n=3+(n-1)\\cdot 2=2n+1$。
(2) $a_{10}=21$，$S_{10}=\\dfrac{10}{2}(a_1+a_{10})=120$。

\par\noindent\textbf{【解析】}使用等差数列通项公式与求和公式。



\item 求函数 $f(x)=x^3-3x$ 的极值。


\par\noindent\textbf{【答案】}$f'(x)=3x^2-3$，令 $f'(x)=0$ 得 $x=\\pm 1$。
$f(1)=-2$ 为极小值，$f(-1)=2$ 为极大值。

\par\noindent\textbf{【解析】}利用导数判定极值。



\item 证明：对任意实数 $a,b$，有
$(a-b)^2 \ge 0$。


\par\noindent\textbf{【答案】}$(a-b)^2\\ge 0$。展开得 $a^2-2ab+b^2\\ge 0$，故命题成立。

\par\noindent\textbf{【解析】}平方恒非负。



\item 求定积分 $\int_0^1 2x\,dx$。


\par\noindent\textbf{【答案】}$\\displaystyle\\int_0^1 2x\\,dx = \\bigl[x^2\\bigr]_0^1 = 1$。

\par\noindent\textbf{【解析】}基本积分公式。


\end{enumerate}
\vspace{1em}


\end{document}